\documentclass[10pt,a4paper]{amsart}
\usepackage[margin=19mm]{geometry}
\usepackage{amsmath,amssymb,mathtools}
\usepackage[scr=rsfs,cal=euler]{mathalfa}
\usepackage{xfrac}
\usepackage[unicode,hidelinks,bookmarks=false]{hyperref}
\newcommand{\ie}{i.e.\ }
\newcommand{\cf}{cf.\ }
\newcommand{\resp}{respectively\ }
\newcommand{\Subsection}{\subsection}
\providecommand{\nobreakdash}{\nobreak\hbox{-}}
\setlength{\emergencystretch}{2em}
\multlinegap0pt
\usepackage{amssymb}
\usepackage{xcolor}
\usepackage{dsfont}
\usepackage[shortlabels]{enumitem}
\usepackage{float}
\newtheorem{theorem}{Theorem}
\def \d {\mathrm{d}}
\title{Maximal regularity for time-fractional Schrödinger equations and application to nonlinear equations}
\author{S. E. Chorfi, F. Et-tahri, L. Maniar, M. Yamamoto}
\date{Source: arXiv:2603.16726v1 (CC BY 4.0)}
\begin{document}
\maketitle

\noindent\emph{Source and licence.} Maximal regularity for time-fractional Schrödinger equations and application to nonlinear equations. Version arXiv:2603.16726v1 (CC BY 4.0), distributed by arXiv under the Creative Commons Attribution 4.0 International licence, http://creativecommons.org/licenses/by/4.0/, as recorded in the arXiv licence metadata for this exact version and shown on the abstract page https://arxiv.org/abs/2603.16726v1. Full licence notice: \texttt{licenses/arxiv-2603.16726-CC-BY-4.0.txt}.\par\smallskip

\noindent\textit{Editorial scope.} Counted unit: the article's theorem pithmMR (source0.tex lines 777--786). The article's complete original proof of it is delivered with it (source0.tex lines 1084--1149, one \texttt{proof} environment of 66 lines, which the article places away from the statement). No mathematics is written, repaired or re-derived: the statement and the proof are the article's own lines, in the article's own notation, and no retained line is substituted or re-flowed.  Retained uncounted as the source-local context the proof needs: lines 80--82; lines 767--769; lines 979--981; lines 1011--1019; lines 1037--1040; lines 1044--1048.  Nothing was omitted from any retained span, so the package carries no omission record.  No retained line contains a citation, so the wrapper carries no bibliography and no bibliography entry.  No retained line contains a figure environment, an image-inclusion command or drawing code, so no image is delivered and the package declares no asset dependency.  Editorial formatting only, in addition: an AMS \texttt{amsart} preamble that restates the article's own declarations and macros used by the retained lines, namely the environments \texttt{theorem} on separate counters, together with the standard packages those lines need.  The retained environments print in this document as Theorem 1 in order of appearance, whereas the article prints the same items with the numbering of its own full document.  The complete native source of the article as distributed in the arXiv e-print for this exact version is bundled unchanged as \texttt{sources/196572/source0.tex}.
\begin{equation}\label{fseq0}
\partial_t^\alpha(u-u_0) -\mathrm{i} Au=f, \qquad t \in (0,T),
\end{equation}

\begin{equation}\label{pifseq}
\partial_t^\alpha (u-u_0) -\mathrm{i} Au=f, \qquad t \in (0,T).
\end{equation}

\begin{equation}\label{semieq0}
\partial_t^\alpha (u(t)-u_0) -\mathrm{i} Au(t)=F(u(t)), \qquad t \in (0,T),
\end{equation}

\begin{theorem}\label{ThmAch}
Let $\alpha\in (0,1)$, $p>\frac{1}{\alpha}$ and assume the assumption \textbf{(H)}. Then, there exists $T_1>0$ such that, for all $u_0\in \mathcal{X}_{\alpha,p}$ and all $T\in (0,T_1)$, the equation \eqref{semieq0} has a unique solution $u\in \mathcal{M}\mathcal{R}_{\alpha,p}(T).$ %Moreover, there exists a constant $C=C(\alpha,p)>0$, independent of $u_0$ and $F$, such that
%\[
%\|u\|_{\mathcal{M}\mathcal{R}_{\alpha,p}(T)}
%\le
%C
%\left(\|u_0\|+\|F(0)\|_{L^{p}(0,T;H)}
%\right).\]
\end{theorem}

\begin{equation}
    k_\alpha^{*n}(t)
= \frac{\Gamma(\alpha)^n}{\Gamma(n\alpha)}\, t^{n\alpha-1}. \label{n_convolution}
\end{equation}

    \begin{align}
		\int_0^T \|u(t)\|^p\d t
		&\le \frac{T^{\frac{\alpha p}{q}}}{\alpha^{\frac{p}{q}}\Gamma(\alpha)^p}\int^T_0 k_{\alpha}(T-r)
		\int_0^r\|\partial_t^\alpha u(s)\|^p\;\d s\;\d r. \label{lm:semi}
	\end{align}

\begin{theorem}\label{pithmMR}
Let $1<p<\infty, \; \alpha >\frac{1}{p}$, $u_0\in \mathcal{X}_{\alpha,p}$ and $f\in L^p(0,T;H).$ Then \eqref{pifseq} satisfies the maximal $L^p$-regularity. That is, the solution given by
\begin{equation*}
u(t)=\sum_{n=1}^{\infty}\left[\left\langle u_0, \varphi_n\right\rangle E_{\alpha, 1}\left(-\mathrm{i} \lambda_n t^\alpha\right) + \int_0^t \langle f(s),\varphi_n\rangle\,(t-s)^{\alpha-1}E_{\alpha,\alpha}\big(-\mathrm{i}\lambda_n (t-s)^\alpha\big)\,\d s \right]\varphi_n,
\end{equation*}
satisfies $u\in L^p(0,T;D(A))$, and there exists a constant $C>0$ such that
\[
\|u-u_0\|_{W_{\alpha,p}(0,T;H)} + \|u\|_{L^p(0,T;D(A))}\le C\left(\|u_0\|_{\mathcal{X}_{\alpha,p}}+\|f\|_{L^p(0,T;H)}\right).
\]
\end{theorem}

\begin{proof}[Proof of Theorem \ref{ThmAch}]
We define the linear operator $\mathcal{L}: L^p(0,T;H)\times\mathcal{X}_{\alpha,p}\to \mathcal{M}\mathcal{R}_{\alpha,p}(T)$ by  $\mathcal{L}(f,u_0)=u,$ where $u$ is the solution to \eqref{fseq0} associated with $u_0$ and $f$. By the maximal regularity (Theorem \ref{pithmMR}), we have $$\mathcal{L}\in \mathcal{B}\left(L^p(0,T;H)\times \mathcal{X}_{\alpha,p}, \mathcal{M}\mathcal{R}_{\alpha,p}(T)\right).$$ 
We fix $u_0\in \mathcal{X}_{\alpha,p}$ and define $\Phi:\mathcal{M}\mathcal{R}_{\alpha,p}(T)\to \mathcal{M}\mathcal{R}_{\alpha,p}(T)$ by
$$\Phi(v):=\mathcal{L}(F(v),u_0).$$
Note that, for all $v\in \mathcal{M}\mathcal{R}_{\alpha,p}(T)$, we have $F(v)\in L^p(0,T;H)$, so $\Phi$ is well defined. Then $u$ solves \eqref{semieq0} if and only if $\Phi(u)=u$. Therefore, we will apply the Banach contraction mapping theorem for iterates of $\Phi$.

Let $v,w\in \mathcal{M}\mathcal{R}_{\alpha,p}(T)$. Then,
 \begin{align*}
     &\|\Phi(v)-\Phi(w)\|_{\mathcal{M}\mathcal{R}_{\alpha,p}(T)}= \|\mathcal{L}(F(v)-F(w),0)\|_{\mathcal{M}\mathcal{R}_{\alpha,p}(T)}\\
     &\le \|\mathcal{L}\| \|F(v)-F(w)\|_{L^p(0,T;H)}\\
     &\le \epsilon \|\mathcal{L}\|\|v-w\|_{\mathcal{M}\mathcal{R}_{\alpha,p}(T)}+C\|\mathcal{L}\|\|v-w\|_{L^p(0,T;H)}\\
     &\le \epsilon \|\mathcal{L}\|\|v-w\|_{\mathcal{M}\mathcal{R}_{\alpha,p}(T)}\\
     &\quad+C\|\mathcal{L}\|\left(\|(v-w)-(v(0)-w(0))\|_{L^p(0,T;H)}+ \|v(0)-w(0)\|_{L^p(0,T;H)}\right)\\
     &= \epsilon \|\mathcal{L}\|\|v-w\|_{\mathcal{M}\mathcal{R}_{\alpha,p}(T)}\\
     &\quad+C\|\mathcal{L}\|\left(\|(v-w)-(v(0)-w(0))\|_{L^p(0,T;H)}+ T^{\frac{1}{p}} \|v(0)-w(0)\|\right).
 \end{align*}
Using the estimate \eqref{lm:semi}, we obtain
\begin{align*}
     &\|\Phi(v)-\Phi(w)\|_{\mathcal{M}\mathcal{R}_{\alpha,p}(T)}\le (\epsilon +CT^{\frac{1}{p}})\|\mathcal{L}\|\|v-w\|_{\mathcal{M}\mathcal{R}_{\alpha,p}(T)}\\
     &+C \|\mathcal{L}\|\frac{T^{\frac{\alpha }{q}}}{\alpha^{\frac{1}{q}}\Gamma(\alpha)^p}\left(\int^T_0 k_{\alpha}(T-r)
		\int_0^r\|\partial_t^\alpha ((v-w)-(v(0)-w(0)))(s)\|^p\;\d s\;\d r\right)^{\frac{1}{p}} \\
    &\le (\epsilon +CT^{\frac{1}{p}})\|\mathcal{L} \|\|v-w\|_{\mathcal{M}\mathcal{R}_{\alpha,p}(T)}
     +C \|\mathcal{L}\|\frac{T^{\frac{\alpha }{q}}}{\alpha^{\frac{1}{q}}\Gamma(\alpha)^p}\left(\int^T_0 k_{\alpha}(T-r)
		\|v-w\|^p_{\mathcal{M}\mathcal{R}(r)}\;\d r\right)^{\frac{1}{p}}\\
        & =\frac{1}{2}\|v-w\|_{\mathcal{M}\mathcal{R}_{\alpha,p}(T)}
     + \gamma\left(\int^T_0 k_{\alpha}(T-r)
		\|v-w\|^p_{\mathcal{M}\mathcal{R}(r)}\;\d r\right)^{\frac{1}{p}},
 \end{align*}
 where we choose $\epsilon=\frac{1}{4\|\mathcal{L}\|}$, $0<T\le \frac{1}{(4 C\|\mathcal{L}\|)^p}$ and set $\gamma:=C \|\mathcal{L}\|\frac{T^{\frac{\alpha }{q}}}{\alpha^{\frac{1}{q}}\Gamma(\alpha)^p}$.
 
Thus, by induction, for every $n\ge 1$, we obtain
\begin{align*}
     \|\Phi^n(v)-\Phi^n(w)\|_{\mathcal{M}\mathcal{R}_{\alpha,p}(T)}
&\le \frac{1}{2^n}\|v-w\|_{\mathcal{M}\mathcal{R}_{\alpha,p}(T)}  \\
&\quad + \gamma
\sum_{j=1}^{n}
\frac{1}{2^{n-j}}
\left(
\int_0^T k_\alpha^{(*j)}(T-r)\,
\|u-v\|_{\mathcal{MR}(r)}^p\,dr
\right)^{\frac{1}{p}}\\
&\hspace{-1cm}\le \left[\frac{1}{2^n}+ \gamma
\sum_{j=1}^{n}
\frac{1}{2^{n-j}}
\left(
\int_0^T k_\alpha^{(*j)}(T-r)\,
\,dr
\right)^{\frac{1}{p}}\right]\|v-w\|_{\mathcal{M}\mathcal{R}_{\alpha,p}(T)}.  
 \end{align*}
 Using the formula \eqref{n_convolution}, we obtain
 \begin{align*}
     \|\Phi^n(v)-\Phi^n(w)\|_{\mathcal{M}\mathcal{R}_{\alpha,p}(T)}
&\le \frac{1}{2^{n}}\left[1+ \gamma
\sum_{j=1}^{n}
2^{j}
\left(
\frac{\Gamma(\alpha)^j}{\Gamma(j\alpha+1)}T^{j\alpha}
\right)^{\frac{1}{p}}\right]\|v-w\|_{\mathcal{M}\mathcal{R}_{\alpha,p}(T)}.
 \end{align*}
The convergence of the series $\displaystyle \sum_{j\ge 1}
2^{j}
\left(
\frac{\Gamma(\alpha)^j}{\Gamma(j\alpha+1)}T^{j\alpha}
\right)^{\frac{1}{p}}$ follows from Stirling's formula and the Cauchy root test. Then, for $n$ sufficiently large, 
$\Phi^n$ is a contraction. By Banach's contraction mapping theorem, there is a unique solution $u\in \mathcal{M}\mathcal{R}_{\alpha,p}(T).$
\end{proof}

\end{document}
